% Generated draft: requires mathematical and source review.
\documentclass[11pt]{article}
\usepackage[margin=25mm]{geometry}
\usepackage{fontspec}
\setmainfont{Latin Modern Roman}
\usepackage{amsmath,amssymb,mathtools,unicode-math}
\setmathfont{Latin Modern Math}
\usepackage{enumitem,needspace,xurl,hyperref}
\setlength{\parindent}{0pt}
\setlength{\parskip}{5pt}
\setlength{\emergencystretch}{3em}
\newcommand{\N}{\mathbb{N}}
\newcommand{\Z}{\mathbb{Z}}
\newcommand{\Q}{\mathbb{Q}}
\newcommand{\R}{\mathbb{R}}
\newcommand{\C}{\mathbb{C}}
\newcommand{\F}{\mathbb{F}}
\newcommand{\A}{\mathbb{A}}
\newcommand{\PP}{\mathbb{P}}
\newcommand{\E}{\mathbb{E}}
\newcommand{\eps}{\varepsilon}
\newcommand{\dd}{\,\mathrm{d}}
\begin{document}
\section*{UnsolvedMath ID: 30006418}
\textbf{Source problem code:} catalogue entry
\subsection{Definitions and mathematical statement}
For metric spaces \((X,d_X)\) and \((Z,d_Z)\), a submetry is a map \(f:X\to Z\) such that
\[
f\bigl(\overline B_X(x,r)\bigr)=\overline B_Z(f(x),r)
\qquad(x\in X,\ r\geq 0),
\]
where \(\overline B_X(x,r)=\{x'\in X:d_X(x,x')\leq r\}\). Two maps \(f:X\to Z\) and \(g:X'\to Z'\) are topologically equivalent if there are homeomorphisms \(h:X\to X'\) and \(k:Z\to Z'\) satisfying \(g\circ h=k\circ f\).

\textbf{Source conjecture, with unresolved hypothesis conventions.} Let \((P_i:M_i\to Y_i)_{i\in\N}\) be a sequence of manifold submetries. Assume that each \(M_i\) is compact, that the manifolds \(M_i\) have uniformly bounded geometry, and that
\[
M_i\xrightarrow{\mathrm{GH}}M.
\]
Assume also that the maps \(P_i\) converge to a submetry \(P:M\to Y\), and that \emph{both} sequences \((M_i)\) and \((Y_i)\) do not collapse. Must there exist \(N\in\N\) such that, for every \(i,j\geq N\), there are homeomorphisms
\[
h_{ij}:M_i\longrightarrow M_j,
\qquad
k_{ij}:Y_i\longrightarrow Y_j
\]
with
\[
P_j\circ h_{ij}=k_{ij}\circ P_i?
\]

The supplied record does not define the precise fiber regularity required by ``manifold submetry,'' the geometric bounds included in ``uniformly bounded geometry,'' the convergence notion for maps between varying spaces, or the meaning of noncollapse for either sequence. These hypotheses are retained as source terminology and require clarification before this becomes a fully specified conjecture. In particular, no particular curvature, dimension, injectivity-radius, volume, or diameter bounds are added here. No manifold structure on the limit \(M\) or \(Y\) is asserted beyond what the supplied hypotheses might imply under their intended conventions.
\subsection{Short English statement}
The conjecture asks whether a convergent sequence of compact manifold submetries eventually has a single topological type as a map, provided the source manifolds have uniformly bounded geometry and neither the source nor target sequence collapses. Equivalence means that homeomorphisms of sources and targets make the corresponding square commute. Several geometric and convergence conventions remain unspecified in the supplied record.
\subsection{Sources}
\begin{itemize}
\item Supplied local UnsolvedMath registry record, permanent ID \(30006418\), title \emph{Stability of Noncollapsing Manifold Submetries}, source problem code \texttt{catalogue entry}; collection: Oberwolfach Reports Open Problems. The supplied record labels the problem open and its statement \texttt{corrected\_verified}; these are supplied curation labels, not independent verification. No proposer is identified; the supplied proposed year is 2025.
\item Supplied report citation: \emph{Differentialgeometrie im Großen (2026)}, Oberwolfach Reports, DOI \texttt{10.4171/owr/2025/30}, pp.~1591--1646. Supplied source URL and literature-source link, labeled ``Oberwolfach Report 30/2025'': \url{https://doi.org/10.4171/owr/2025/30}. The supplied report title year and report-number/DOI year differ and have not been reconciled.
\item The supplied report excerpt identifies the statement as Conjecture~1. Its surrounding discussion reports validity for one-dimensional base spaces and cites Perelman's stability theorem as reference [2], Harvey's equivariant stability theorem as reference [1], and Tapp's result for Riemannian submersions as reference [3]. Full bibliographic entries and the hypotheses of those results are not supplied; these statements are retained as attributed context.
\item Literature assessment dated exactly \texttt{2026-08-22}: the supplied curation labels the conjecture open, reports only special cases known, and says that no general proof or counterexample was found. Its listed source checked is Oberwolfach Report 30/2025 at \url{https://doi.org/10.4171/owr/2025/30}. This is supplied dated curation, not an independently verified literature assessment; the record itself requests independent verification of time-sensitive status and cited claims.
\end{itemize}
\end{document}
